\section{Introduction}

The Cremona group $\Cr_2(k)$ over a field $k$ is the group of birational transformations of the projective plane $\p^2$ over $k$. Cremona groups have been the delight of algebraic geometers and group theorists in both, classical and modern times. As of today, many aspects of $\Cr_2(k)$ are well understood and there are many tools at hand to study those groups. For instance, $\Cr_2(k)$ acts by isometries on an infinite dimensional hyperbolic space $\mathbb{H}^\infty$ \cite{Cantat_groupes_birat} and on various CAT(0) cube complexes \cite{lonjou2021actions}. Nevertheless, some questions have remained open. The goal of this article is to positively answer a question asked more than a decade ago by Favre in \cite[Quest.\,1]{favrebourbaki}.

Let us fix projective coordinates $[x:y:z]$ of $\p^2$. An element $f\in\Cr_2(k)$ is given~by
\[
[x:y:z]\dashmapsto [f_0:f_1:f_2],
\]
where the $f_i\in k[x,y,z]$ are homogeneous of the same degree and without non-constant common factor. The \emph{degree} $\deg(f)$ of $f$ is defined as the degree of the polynomials~$f_i$. An element $f\in\Cr_2(k)$ is \emph{algebraic}, if $\deg(f^n)$ is uniformly bounded for all $n\in\Z$. A subgroup $G<\Cr_2(k)$ is \emph{bounded} if the degree of all elements in $G$ is uniformly bounded.
Clearly, a bounded subgroup consists of algebraic elements. However, the converse is not true. For instance, consider the subgroup defined in the affine coordinates $(x,y)$ of $\p^2$ by
\[
G=\{(x,y)\dashmapsto (x+R(y), y)\mid R\in k(y)\}.
\]
Every element in $G$ is algebraic, but $G$ is not bounded. In this paper, we show the following theorem, which solves Favre's question:

\begin{theorem}\label{thm:main}
{Let $k$ be a field and} let $G<\Cr_2(k)$ be a finitely generated subgroup such that every element of $G$ is algebraic. Then $G$ is bounded.
\end{theorem}

\begin{remark}\label{rmk_alg_closed_fields}
Note that the properties of being algebraic and being bounded are invariant under field extensions, so it is enough to show Theorem~\ref{thm:main} for algebraically closed fields.
\end{remark}

\begin{remark}\label{rmk_redu_Cantat}
The first step towards the proof of Theorem~\ref{thm:main} is due to Cantat, who showed in \cite{Cantat_groupes_birat} that a finitely generated subgroup $G<\Cr_2(k)$ consisting of algebraic elements is bounded or preserves a rational fibration (see \cite{Lamy_Cremona} for a proof of this result for fields of arbitrary characteristic). It is therefore enough to show Theorem~\ref{thm:main} for finitely generated groups that preserve a rational fibration.
\end{remark}

{In order to show Theorem~\ref{thm:main} for finitely generated groups that preserve a rational fibration}, we introduce the \emph{Jonquières complex} --- a $\mathrm{CAT}(0)$ cube complex with an isometric action of the group of birational transformations preserving a rational fibra\-tion, whose vertex stabilizers are bounded subgroups. Using a suitable description of this complex and the dynamics of $\PGL_2(k)$ on $\p^1$, we show that the action of our group on this complex is decent, as defined below.

\begin{definition}
\label{def:decent}
Let $X_0$ be a set. We say that a group $G_0$ acts on $X_0$ \emph{purely elliptically}, if each element of $G_0$ fixes a point of $X_0$. We say that a group $G_0$ acts on~$X_0$ \emph{decently} if
\begin{itemize}
\item
each subgroup of $G_0$ with a finite orbit fixes a point of $X_0$, and
\item
each finitely generated subgroup of $G_0$ acting purely elliptically fixes a point of~$X_0$.
\end{itemize}
\end{definition}

It is an easy exercise
that if $G_0$ is the isometry group of a simplicial tree $X_0$, then~$G_0$ acts on $X_0$ decently {(see \cite[Cor.\,3 in \S6.5 \& Ex.\,6.3.4]{Serre})}. More generally, if~$G_0$ is the isometry group of a $\mathrm{CAT}(0)$ cube complex $X_0$ with no infinite cubes, then~$G_0$ acts on $X_0$ decently \cite{GLU_Neretin}. Similarly, if~$G_0$ is the isometry group of a $\mathrm{CAT}(0)$ $2$-complex $X_0$ with rational angles, then $G_0$ acts on~$X_0$ decently \cite{NOP}. For further examples, see \cite{HO}.

\subsection{Applications.}
The Cremona group $\Cr_2(k)$ can be equipped with the \emph{Zariski topology} (see for instance \cite{Serre_Bourbaki}). An \emph{algebraic subgroup} of $\Cr_2(k)$ is a Zariski closed bounded subgroup. This explains the terminology: an element in $\Cr_2(k)$ is algebraic if and only if it is contained in an algebraic subgroup. Note that for any $d\geq1$, the subset of $\Cr_2(k)$ consisting of elements of degree at most $d$ is closed (see \cite{Blanc_Furter}), hence a bounded subgroup is contained in an algebraic subgroup. An algebraic subgroup $G$ is always \emph{projectively regularizable}, \ie there exists a birational map $\varphi\colon \p^2\dashrightarrow S$ such that $\varphi G\varphi^{-1}<\Aut(S)$ for some regular projective surface $S$. This follows from the theorems of Weil and Sumihiro or from the fact that the number of base-points of elements in an algebraic subgroup is uniformly bounded (we refer to \cite{Lamy_Cremona} for references and a proof of this fact, or to \cite{lonjou2021actions}). Theorem~\ref{thm:main} has therefore the following direct consequence:

\begin{corollary}
\label{cor:main}
{Let $k$ be a field and} let $G<\Cr_2(k)$ be a finitely generated subgroup such that every element of $G$ is algebraic. Then $G$ is projectively regularizable. \end{corollary}

From another point of view, algebraic elements correspond exactly to elements in~$\Cr_2(k)$ inducing an elliptic isometry on the infinite dimensional hyperbolic space~$\mathbb{H}^\infty$, on which $\Cr_2(k)$ acts \cite{Cantat_groupes_birat}. Theorem~\ref{thm:main} therefore states that the action of $\Cr_2(k)$ on~$\mathbb{H}^\infty$ is decent.

While all algebraic elements are projectively regularizable, there exist also non-algebraic elements that are projectively regularizable. It is still unknown whether a finitely generated subgroup of the Cremona group containing only projectively regularizable elements is projectively regularizable or not. This question is equivalent to the question whether $\Cr_2(k)$ acts decently on the \emph{blow-up complex} --- a CAT(0) cube complex constructed in \cite{lonjou2021actions}. In \cite{GLU_Neretin}, the first and third authors together with Genevois positively answer this question when the base-field {$k$} is finite.

Since elements in $\Cr_2(k)$ preserving a rational fibration are projectively regularizable if and only if they are algebraic {(see for instance \cite[Table\,1]{lonjou2021actions})}, Corollary~\ref{cor:main} directly implies the following.
Let $G<\Cr_2(k)$ be a finitely generated subgroup preserving a rational fibration such that every element in $G$ is projectively regularizable. Then $G$ is projectively regularizable.

The notions of algebraic elements {and bounded subgroups} generalize to birational transformations of arbitrary regular projective surfaces and Theorem~\ref{thm:main} generalizes to this setting:

\begin{theorem}\label{thm:main_general}
Let $S$ be a regular projective surface over a field $k$ and let $G <\Bir(S)$ be a finitely generated subgroup such that every element of~$G$ is algebraic. Then $G$ is {bounded}.
\end{theorem}

However, the most interesting and difficult case is the one of rational surfaces. In~order to keep the notation more accessible, we discuss and prove the general case in the separate Section~\ref{sec:remarks}.

In Section~\ref{sec:degrees} we apply Theorem~\ref{thm:main} to give a first description of the asymptotic degree growth of finitely generated subgroups of $\Cr_2(k)$. This opens up new interesting questions about the dynamical behaviour of finitely generated subgroups of~$\Cr_2(k)$. Let $G<\Cr_2(k)$ be a finitely generated subgroup with a finite generating set $T$. Denote by $B_T(n)$ the set of all elements in $G$ of word length $l_T$ at most $n$ with respect to the generating set $T$, and define
\[
D_{T}(n):=\max_{f\in B_T(n)}\{\deg(f)\}.
\]
It has been shown in \cite{MR3818624} that there are only countably many integer sequences that can appear in this way. However, still very little is known about them. In the case where $T$ consists of a single element, the growth of the function $D_T$ has been {extensively} studied (see Theorem~\ref{thm:hyperboloid}). Theorem~\ref{thm:main} is the main ingredient for the following result, which we prove in Section~\ref{sec:degrees}. For two functions, $f$ and $g$ on $\N$, we write $f\asymp g$ if $f(x)\leq ag(bx)$ and $g(x)\leq cf(dx)$ for some $a,b,c,d>0$.
\begin{corollary}\label{cor:degree}
Let $\kk$ be an algebraically closed field and let $G<\Cr_2(k)$ be a finitely generated subgroup with generating set $T$. Then one of the following is true:
\begin{enumerate}
\item All elements in $G$ are algebraic, $G$ fixes a point in $\h$, and $D_T(n)$ is bounded.
\item[(2a)] The group $G$ contains an element that induces a parabolic isometry of $\h$, $D_T(n)\asymp n$, and $G$ preserves a rational fibration.
\item[(2b)] The group $G$ contains an element that induces a parabolic isometry of $\h$, $D_T(n)\asymp n^2$, and $G$ preserves an elliptic fibration.
\item[(3)] The group $G$ contains an element that induces a loxodromic isometry of $\h$ and $D_T(n)\asymp\lambda^n$ for $\lambda>0$.
\end{enumerate}
\end{corollary}

It would be interesting to study the dynamical behaviour of the degrees of finitely generated subgroups in more detail. In Corollary~\ref{cor:degree}, the asymptotic behaviour does not depend on the choice of the finite generating set $T$, {since for any other generating set $T'$ there is a constant $d>0$ such that for all $n$ we have $B_{T'}(n)\subseteq B_T(dn)$, and so $D_{T'}(n)\leq D_{T}(dn)$.} For instance, one could ask, after fixing $T$, what is the precise asymptotic growth of $D_T(n)$, if the group $G$ contains a loxodromic element:

\begin{question}
Let $G<\Cr_2(k)$ be a finitely generated subgroup containing an element that induces a loxodromic isometry of $\h$, with finite generating set $T$. What is the asymptotic growth of $D_T(n)$? Do we always have $D_T(n)\sim c\lambda^n$ for some constants $c,\lambda>0$? Here, we write $f\sim g$ if $f$ and $g$ are functions such that $\lim_{n\to\infty} \sfrac{f(n)}{g(n)}=1$.

\end{question}

In another direction, it has been shown that the \emph{dynamical degree} of $f\in\Cr_2(k)$, \ie the number $\lambda(f):=\lim_{n\to\infty}\deg(f)^{1/n}$ is always an algebraic integer (more precisely, it is always a Pisot or Salem number, see \cite{blanc2016dynamical}). This leads to the following natural question, which is due to Cantat:

\begin{question}
Let $G<\Cr_2(k)$ be a finitely generated subgroup with finite generating set $T$. Which real numbers can be realized as $\lim_{n\to\infty}D_T(G)^{1/n}$?
\end{question}

For instance, it could be interesting to construct examples of such subgroups and generating sets such that $\lim_{n\to\infty}D_T(G)^{1/n}$ is transcendental. If $S$ is a {regular} projective surface and $f\in\Aut(S)$, then $\lambda(f)$ is the spectral radius of the induced transformation $f_*$ of $f$ on the Neron--Severi lattice of $S$. For this reason, the limit $\lim_{n\to\infty}D_T(G)^{1/n}$ should be seen as an analogue to the \emph{joint spectral radius}, which has been introduced in \cite{rota1960note} and has been studied by many authors (see for instance \cite{jungers2009joint} or \cite{breuillard2021joint}).

\subsection{Decent actions.}
In Section~\ref{sec:Jon} we will show that Theorem~\ref{thm:main} is a special case of a following more general result on groups acting decently on an infinite restricted product.

\label{sum}
A \emph{pointed set} $(X_0,x_0)$ is a set $X_0$ and a point $x_0\in X_0$. The \emph{restricted product} $\bigoplus_{p\in P} (X_p,x_p)$ of a family $\{(X_p,x_p)\}_{p\in P}$ of pointed sets is the set of sections $\{y_p\}_{p\in P}$ with $y_p\in X_p$ such that all but finitely many $y_p$ are equal to $x_p$.

Note that for finite $P$ we have $\bigoplus_{p\in P} (X_p,x_p)=\Pi_{p\in P} X_p$. For infinite $P$, we have that
$\bigoplus_{p\in P} (X_p,x_p)$ is a proper subset of $\Pi_{p\in P} X_p$.

If each $X_p$ is a simplicial tree and each $x_p$ is a vertex (which will be the case in our application towards Theorem~\ref{thm:main}), then $\bigoplus_{p\in P} (X_p,x_p)$ has a structure of a cube complex {whose cubes have the form $\Pi_{p\in P} I_p$, where all but finitely many $I_p$ are equal to $x_p$ and the remaining $I_p$ are edges of $X_p$}. This cube complex is $\CAT(0)$ though this will not be exploited explicitly in the current article.

Let $G_0$ be a group acting on $X_0$ and let $H$ be a group acting on~$P$. The (unrestricted) \emph{wreath product} $G_0\wr_P H$ of $G_0$ and~$H$ over $P$ is the semi-direct product of $\Pi_{p\in P} G_p$, where $G_p=G_0$, with the group~$H$ acting on $\Pi_{p\in P} G_p$ by $h\cdot \{g_p\}_{p\in P}=\{g_{h^{-1}(p)}\}_{p\in P}$.

Let $x_0\in X_0$. We will be considering the subgroup $G^\oplus$ of $G_0\wr_P H$ preserving $\bigoplus_{p\in P} (X_p,x_p)$, where $X_p=X_0$ and $x_p=x_0$, and where the action is defined as follows. For $g=\{g_p\}_{p\in P}\in \Pi_{p\in P}G_p$, we have $g \cdot \{y_p\}_{p\in P}=\{g_p(y_p)\}_{p\in P}$. For $h\in H$, we have $h\cdot \{y_p\}_{p\in P}=\{y_{h^{-1}(p)}\}_{p\in P}$.

It is not hard to verify that if $X_0$ is a simplicial tree, and $f$ is a combinatorial isometry of the cube complex $\bigoplus_{p\in P} (X_p,x_p)$, then $f$ induces a bijection of~$P$. Furthermore, if $G_0$ is the group of all simplicial isometries of $X_0$, then we have $f\in G^\oplus$ if and only if $H$ contains that bijection. For example, if $P=\{p,q\}$ with $X_0$ the real line tiled by unit intervals, then $\bigoplus_{p\in P} (X_p,x_p)$ is the square tiling of the Euclidean plane, and a $90^\circ$ rotation of the plane induces the bijection interchanging $p$ and $q$.

\begin{theorem}
\label{thm:maing}
Let $G_0$ be a group acting decently on $X_0$. Let $P$ be the set of {the points of the projective line $\p^1(k)$} over an algebraically closed field $k$, and let $H=\Aut(\p^1)$. Then $G^\oplus$ acts decently on $\bigoplus_{p\in P} (X_p,x_p)$. \end{theorem}

We will prove Theorem~\ref{thm:maing} in Section~\ref{sec:main}. {We will deduce Theorem~\ref{thm:main} from Theorem~\ref{thm:maing} at the end of Section~\ref{sec:Jon}}.
Note that in Theorem~\ref{thm:maing}, we have to make some assumptions on $P$ and $H$. For example, assume that the set $X_0$ is {a finite simplicial tree and $G_0$ is a nontrivial subgroup of the group of the simplicial isometries of~$X_0$}. Suppose that $H$ is a finitely generated infinite torsion group acting by left multiplication on $P=H$. Then the entire $G^\oplus$ is finitely generated and acts purely elliptically on $\bigoplus_{p\in P} (X_p,x_p)$, but does not have a fixed-point.

\subsubsection*{Acknowledgements}
We thank Serge Cantat for his valuable comments that helped to improve the exposition of this article. {We also thank Chris Karpinski and the anonymous referees for their helpful suggestions.}
The first author would like to thank the CRM (Centre de Recherche Mathématiques in Montreal), the Simons foundation and the organizers of the thematic semester ``Théorie géométrique des groupes'' for her stay at the CRM where a part of this project was realized.

\section{Preliminaries}
In this section, we briefly recall some well-known facts about blow-ups and conic bundles. {By Remark \ref{rmk_alg_closed_fields}, it is enough to work over an algebraically closed field for our problem, hence} we assume our base-field $k$ to be algebraically closed. Unless mentioned otherwise, surfaces are assumed to be projective and smooth.

\subsection{Subgroups of algebraic elements of automorphism groups}
{Let $S$ be a surface over $k$.} An ample divisor $H$ on $S$ defines a degree function $\deg_H$ on $\Bir(S)$ by $\deg_H(f)=f_*H\cdot H$. An element $g\in \Bir(S)$ is then \emph{algebraic} if the degree sequence $\{\deg_H(f^n)\}_n$ is bounded. A subgroup $G<\Bir(S)$ is \emph{bounded} if $\{\deg_H(f)\mid f\in G\}$ is bounded. The properties of being algebraic and being bounded do not depend on the choice of an ample divisor (see for example \cite{MR4133708}).

Let us observe that Theorem~\ref{thm:main} holds if we work with finitely generated groups of automorphism groups:

\begin{lemma}\label{lem:aut}
Let $S$ be a
surface and let $G<\Aut(S)$ be a finitely generated subgroup such that every element in $G$ is algebraic. Then $G$ is bounded.
\end{lemma}

\begin{proof}
The linear action of $\Aut(S)$ on the Néron--Severi lattice $N^1(S)$ of $S$ yields a homomorphism $\Aut(S)\to\GL(N^1(S))\simeq \GL_n(\Z)$, whose kernel is an algebraic group and therefore bounded. The image of an element $g\in G$ in $\GL(N^1(S))$ is of finite order, since $g$ is algebraic (see for example \cite[Th.\,4.6]{cantat2015cremona}). Using that $G$ is finitely generated, we obtain that the image of $G$ in $\GL(N^1(S))$ is finite. Therefore, $G$ is a finite extension of a bounded group and therefore bounded itself.
\end{proof}

\subsection{Bubble space, strong factorization, and base-points}
Let $S$ be a surface and let $s\in S$. Then there exists a surface $\tilde S$ and a morphism $\pi\colon\tilde S\to S$ such that the fiber $E$ over $s$ is isomorphic to $\p^1$ and $\pi$ induces an isomorphism between $\tilde S\setminus E$ and $S\setminus\{s\}$. The morphism $\pi\colon \tilde S\to S$ is called the \emph{blow-up} of $S$ in $s$, and it is unique up to isomorphism. Consider two distinct points $s,s'$ on a surface $S$ and their respective blow-ups $\pi_{s}$ and~$\pi_{s'}$. Blowing-up first~$s$ and then~$s'$ gives the same transformation up to isomorphism as blowing-up first~$s'$ and then~$s$. By abuse of language, we will say that $\pi_s$ and~$\pi_{s'}$ \emph{commute} and we denote the successive blow-up of~$s$ and then~$s'$ by $\pi_{s'}\pi_s=\pi_s\pi_{s'}$.

The \emph{bubble space} of $S$ is the set $\B(S)$ of triples $(t,T,\pi)$, where $\pi\colon T\to S$ is a birational morphism from a surface $T$ and $t$ is a point on $T$; two triples $(t,T,\pi)$ and $(t', T',\pi')$ are identified if $\pi^{-1}\pi'$ is a local isomorphism around~$t'$ mapping~$t'$ to~$t$. The bubble space can be thought of as the set of all points on $S$ and on surfaces obtained from~$S$ by successively blowing up points. The points in $\B(S)$ contained in~$S$ are called \emph{proper} points.

Zariski's strong factorization theorem states that every birational transformation $f\colon S\dashrightarrow T$ between surfaces can be factored into blow-ups of points. More precisely, there exists a surface $Z$ and a factorization
\[
\begin{tikzcd}
& Z \arrow[ld, "\ts\pi"'] \arrow[rd, "\ts\rho"] & \\
S \arrow[rr, "\ts f", dashed] & & T
\end{tikzcd}
\]
where $\pi$ and $\rho$ are compositions of blow-ups of points. Note that the points blown up by $\pi$ and $\rho$ can be seen as elements of the bubble space $\B( S)$ and $\B( T)$ respectively. Moreover, $Z$ can be chosen in such a way, and will be denoted by $Z_f$, that for any other such factorization
\[
\begin{tikzcd}
& Z' \arrow[ld, "\ts\pi'"'] \arrow[rd, "\ts\rho'"] & \\
S \arrow[rr, "\ts f", dashed] & & T
\end{tikzcd}
\]
there exists a surjective morphism $\eta\colon Z'\to Z_f$ such that $\pi'=\pi\eta$ and $\rho'=\rho\eta$. If we require $Z_f$ to be minimal in this sense, this factorization is unique up to isomorphism and up to possibly changing the order of blowing up the points. {The morphism $\pi\colon Z_f\to S$ is called the \emph{minimal resolution} of $f$}, and the points (in~$\B(S)$) blown up by $\pi$ are called the \emph{base-points} of $f$.

A base-point $s$ of $f$ is called \emph{persistent} if there exists $l\geq 1$ such that, for all $n\geq l$, the point $s$ is a base-point of $f^n$ but $s$ is not a base-point of $f^{-n}$.

\subsection{Conic bundles}

A \emph{rational fibration} on a surface $S$ is a morphism $\pi\colon S\to C$, where $C$ is a
curve, such that all the fibers are isomorphic to $\P^1$. Note that this is a more restrictive definition than the more usual one, which only asks that the fibers are rational curves.

Let $\pi\colon S\to C$ and $\pi'\colon S'\to C$ be rational fibrations.
We say that a birational transformation $f \colon S\dashrightarrow S'$ \emph{preserves} the fibrations if there exists an automorphism $\hbar(f)\colon C\to C$ such that the following diagram commutes
\begin{center}
\begin{tikzcd}
S \arrow[d, "\ts\pi",swap] \arrow[rr, "\ts f", dashed] & & S' \arrow[d, "\ts\pi'"] \\
C \arrow[rr, "\ts\hbar(f)"] & & C.
\end{tikzcd}
\end{center}

{ Let $\pi\colon S\to C$ be a rational fibration. An \emph{elementary transformation} $f\colon S\to S'$ is the composition of blowing up a point $s\in S$ followed by blowing down the strict transform of the fiber containing $s$. The surface $S'$ comes equipped with the rational fibration $\pi'\colon S'\to C$ defined by {$\pi'=\pi f^{-1}$}. Note that $f$ preserves these fibrations for $\hbar(f)=\id$.}

The following fact is well-known (see for instance \cite[Cor.\,3.2]{schneider2022relations}):
\begin{proposition}\label{prop:elementary}
Let $\pi\colon S\to C$ and $\pi'\colon S'\to C$ be rational fibrations.
Every birational transformation $f\colon S\dashrightarrow S'$ preserving the fibrations can be factored { into a sequence of elementary transformations and a (fibration preserving) isomorphism}.
\end{proposition}

A \emph{conic bundle} is a composition with a rational fibration $\pi\colon S\to C$ of a sequence of blow-ups $\tilde S\to S$, such that we blow up {in total} at most one point~$s\in S$ in each fiber. {In other words, $\tilde S\to C$ has only finitely many singular fibers, each of them consisting of
two rational curves of self-intersection $-1$, called also \emph{$-1$-curves} (the preimage of~$s$ under the blow-up and the strict transform of the fiber of $s$).}

\section{The Jonquières complex}
\label{sec:Jon}
In this section we reduce Theorem~\ref{thm:main} to Theorem~\ref{thm:maing}. We always assume our base-field $k$ to be algebraically closed, {which is justified by Remark \ref{rmk_alg_closed_fields}, and we assume the} surfaces to be projective and smooth.

\subsection{The blow-up complex}
Let $T$ be a
surface. In \cite{lonjou2021actions}, the authors constructed the \emph{blow-up complex} $\CC(T)$ --- a $\CAT(0)$ cube complex with an isometric action of $\Bir(T)$. Let us briefly recall the construction of $\CC(T)$. The vertices of $\CC(T)$ are equivalence classes of \emph{marked surfaces}, \ie pairs $(S,\varphi)$, where $S$ is a
surface and $\varphi\colon S\dashrightarrow T$ a birational transformation. Marked surfaces $(S,\varphi)$ and $(S',\varphi')$ are \emph{equivalent} if $\varphi^{-1}\varphi'\colon S'\to S$ is an isomorphism.
Vertices (represented by) $(S,\varphi)$ and $(\tilde{S},\tilde{\varphi})$, are connected by an edge if ${\varphi}^{-1}\tilde \varphi\colon \tilde S\to {S}$ is the blow-up of a point (or its inverse). More generally, $2^n$ different vertices form an $n$-cube if there is a marked surface $(S,\varphi)$ and distinct points $s_1,\dots, s_n\in S$ such that these vertices are obtained by blowing up subsets of $\{s_i\}_{\{1\leq i \leq n\}}$.

\subsection{The Jonquières complex}
Let $\F_0=\p^1\times\p^1$ with the rational fibration onto the first factor $\pi_0\colon \F_0\to\p^1$. We define the \emph{Jonquières complex} $\mathscr{J}$ as the subcomplex of the blow-up complex $\mathcal C(\F_0)$ induced on the set of vertices represented by marked surfaces $(S,\varphi)$ such that, for $\pi=\pi_0\varphi$, the rational map
$\pi\colon S\dashrightarrow\p^1$ is a conic bundle.

\begin{remark}\label{remark_cube_Jonquièrescplx}
Consider $2^n$ vertices of $\mathscr{J}$ spanning a cube. This means that there exists a surface $(S,\varphi)$ and distinct points $s_1,\dots, s_n\in S$ such that these vertices are obtained by blowing up subsets of $\{s_i\}_{\{1\leq i \leq n\}}$. Since these vertices belong to the Jonquières complex, we have that the points~$s_i$ belong to distinct fibers of $\pi$. In other words, we have $\pi(s_i)\neq\pi(s_j)$ for all $i\neq j$.
\end{remark}

\begin{remark}\label{remark:not_convex}
Note that the Jonquières complex $\Jon$ is not a convex subcomplex of the blow-up complex $\mathcal C(\F_0)$. Indeed, consider the birational transformation \hbox{$\varphi:(x,y)\mapsto (x,x^2+y)$}. It has three base-points, but only one of them is proper. Let $\rho \colon S\to \F_0$ be the minimal resolution of $\varphi$. Then the vertex represented by~$(S,\rho)$ is not a vertex of the Jonquières complex because more than one point has been blown-up in the same fiber. However, the vertex represented by $(S,
\rho)$ lies on a geodesic edge-path between the vertices $(\F_0,\id)$ and $(\F_0, \varphi^{-1})$.

Nevertheless, by Proposition~\ref{prop:elementary}, the $1$-skeleton of $\mathscr{J}$ it is isometrically embedded in the $1$-skeleton of $\mathcal C(\F_0)$.
\end{remark}

For any $p\in \p^1$, consider the subcomplex $X_p$ of $\mathscr{J}$ induced on the vertices represented by the marked surfaces $(S,\varphi)$, where $\varphi\colon S\dashrightarrow \F_0$ induces an isomorphism between $S\setminus \pi^{-1}(p)$ and $\F_0\setminus \pi_0^{-1}(p)$.

\begin{lemma}
\label{lem:tree}
For any $p\in \p^1$, the subcomplex $X_p$ is a tree.
\end{lemma}

\begin{proof}
First note that $X_p$ is connected as a consequence of Proposition \ref{prop:elementary} and the fact that two elementary transformations performed in distinct fibers commute.

Second, consider the vertex $v_0=(\F_0,\id)$ in $X_p$, and let $\p^1_p\subset \F_0$ be the fiber over~$p$. Let $v_0,v_1,v_2,v_3,\ldots$ be the consecutive vertices on an edge-path without backtracks in $X_p$. The surface $v_1$ is obtained from $v_0$ by blowing up a point {$s$} in $\p^1_p$ to a $-1$-curve. The surface $v_2=(S_2,\varphi_2)$ is obtained from $v_1$ by blowing down the other $-1$-curve in the fiber over $p$ to a point $s_2$. In particular, the birational transformation $\varphi^{-1}_2$ sends the entire {$\p^1_p\setminus \{s\}$} to $s_2$.

Continuing, for $i\geq 2$ the surface $v_{2i-1}$ is obtained from $v_{2i-2}$ by blowing up a point in the fiber over $p$ distinct from $s_{2i-2}$ to a $-1$-curve, and the surface $v_{2i}$ is obtained from $v_{2i-1}$ by blowing down the other $-1$-curve in the fiber over $p$ to a point $s_{2i}$. Thus the birational transformation $\varphi^{-1}_{2i}$ sends the entire {$\p^1_p\setminus \{s\}$} to $s_{2i}$. In particular, $\varphi_{2i}$ {is not an isomorphism}, and so $v_{2i}\neq v_0$. This proves that there is no cycle in $X_p$, and so $X_p$ is a tree.
\end{proof}

We choose a preferred vertex
$x_p=(\F_0,\id)$ in each $X_p$, and we consider the family of pointed sets $\{(X_p,x_p)\}_{p\in \p^1}$.

\begin{lemma}\label{lem_product_tress}
The cube complex $\mathscr{J}$ is isomorphic to $\bigoplus_{p\in \p^1} (X_p,x_p)$.
\end{lemma}

\begin{proof}[Proof of Lemma~\ref{lem_product_tress}]
Let $(S,\varphi)$ be
a vertex of $\mathscr{J}$. The birational transformation $\varphi \colon S\dashrightarrow\! \F_0$ is decomposed as $\varphi=\sigma_n\cdots\sigma_1$, where each $\sigma_i$ is the blow-up of a point or the blow-down of a $-1$-curve in the fiber over a point $p_i\in \p^1$. For $p_i\neq p_j$,
we have that $\sigma_i$ and $\sigma_j$ commute. Thus there is a finite subset $Q\subset \p^1$ such that
$\varphi=\Pi_{p\in Q}\varphi_p$, where each (of the commuting) $\varphi_p$ is a composition of~$\sigma_i$ with $p_i=p$.
Consider the marked surfaces $\varphi_p\colon S_p \dashrightarrow\F_0$ for $p\in Q$, and $\varphi_p=\id, S_p=\F_0$ for $p\in \p^1\setminus Q$. They represent vertices of~$X_p$. Then $\{(S_p,\varphi_p)\}_{p\in \p^1}$
defines a vertex of $\bigoplus_{p\in \p^1} (X_p,x_p)$.

This is a bijective correspondence between the vertices of $\mathscr{J}$ and $\bigoplus_{p\in \p^1} (X_p,x_p)$, and it extends to an isomorphism on the entire complexes.
\end{proof}

The coordinate $y_p$ of a point $y\in \bigoplus_{p\in \p^1} (X_p,x_p)=\mathscr{J}$ should be thought of as the ``marked fiber''
in the surface corresponding to $y$ over the point $p$. Thus modifying the coordinate $y_p$ of $y$ corresponds to performing an alternating sequence of blow-ups of points and blow-downs of $-1$ curves in the fiber over $p$ of the surface corresponding to $y$.

\smallskip

The \emph{Jonquières group} is the subgroup of $\Bir(\F_0)$ consisting of the \emph{Jonquières transformations} $f$ that preserve the rational fibration $\pi_0\colon \F_0\to \p^1$. {The Jonquières group acts on the vertex set of $\Jon$ by $f\cdot(S, \varphi)=(S, f\varphi)$,
and this action extends to an action by isometries on the entire $\Jon$.}

\begin{remark}\label{rmk_subjonq_fixpoint}
A subgroup $G$ of the Jonquières group is a subgroup of $\Aut(S)$ for some
conic bundle $\pi \colon S\to \p^1$ if and only if $G$ fixes a point in $\Jon$. This is because if $G$ fixes an interior point of a cube in~$\mathscr{J}$ described via a surface $S$ in Remark~\ref{remark_cube_Jonquièrescplx}, then $G$ {fixes this cube and in particular $G$} fixes the vertex corresponding to~$S$. {By~the definition of a marked surface this is equivalent to} $G< \Aut(S)$.
\end{remark}

{We conclude with the following.}

\begin{proof}[Proof of Theorem~\ref{thm:main}]
By Remark \ref{rmk_alg_closed_fields}, we can assume the {base-field $k$} to be algebraically closed. By Remark \ref{rmk_redu_Cantat}, we can assume $G$ to be a subgroup of the Jonquières group.

As a consequence of Lemma \ref{lem_product_tress}, the complex $\bigoplus_{p\in \p^1} (X_p,x_p)$ inherits the action of the Jonquières group from $\Jon$. Because $X_p$ are trees (see Lemma~\ref{lem:tree}), {their isometry groups are decent}. All the elements of the Jonquières group induce elements of $\Aut(\p^1)$ on $\p^1$.
Moreover, by Remark \ref{rmk_subjonq_fixpoint} and by Lemma~\ref{lem:aut}, to prove Theorem~\ref{thm:main} for a subgroup $G$ of the Jonquières group, we need to find a fixed-point for $G$ in $\Jon$, {which is guaranteed by} Theorem~\ref{thm:maing}.
\end{proof}

\section{Proof of the main theorem}
\label{sec:main}

In this section we prove Theorem~\ref{thm:maing}. We keep the notation $\Jon$ for the general $\bigoplus_{p\in \p^1} (X_p,x_p)$ (and not just for the Jonquières complex). We denote by $\hbar$ the quotient map $G_0\wr_P H\to H$. Abusing the notation, for $f\in G_0\wr_P H$ and $p\in P$, we denote by $f(p)\in P$ the point $\hbar(f)(p)$.

\begin{conv}
\label{conv:fibre}
Note that for $f=hg\in G^\oplus$, with $h\in H$ and $g\in\Pi_{p\in P} G_p$, and for $y\in \Jon$, the value $f(y)_p$ equals $g_{h^{-1}(p)}\big(y_{h^{-1}(p)}\big)$ and so it depends only on $h,g_{h^{-1}(p)}$ and $y_{h^{-1}(p)}$. Henceforth, slightly abusing the notation, we will refer to $f(y)_p$ as $f(y_{f^{-1}(p)})$.
\end{conv}

This conveys the fact that for a Jonquières transformation $f$, the ``marked fiber'' in $f(S)$ over $p$ depends only on $f$ and the ``marked fiber'' in $S$ over $f^{-1}(p)$.

{
\begin{lemma}
\label{lem:first_item}
Let $G_0$ be a group acting on $X_0$ so that each subgroup of $G_0$ with a finite orbit fixes a point of $X_0$ (which is the first item of Definition~\ref{def:decent} for $G_0$). Then each subgroup of $G<G^\oplus$ with a finite orbit $Y$ in $\Jon$ fixes a point of $\Jon$ (which is the first item of Definition~\ref{def:decent} for $G$).
\end{lemma}}

\begin{proof}
{
For $p\in P$, let $Y_p=\{y_p \colon y\in Y\}$ denote the finite set of the coordinates of~$Y$ in the factor $X_p$.}

{First, we consider each finite orbit $O\subset P$ of $\hbar(G)<H$. If for all $p\in O$ we have $Y_p=\{x_p\}$, then we define $z_p=x_p$ for all $p\in O$.
Since $Y\subset \Jon$, there are only finitely many finite orbits $O$ of $\hbar(G)<H$ with $Y_p\neq \{x_p\}$ for some $p\in O$. For each such $O$, we choose $o\in O$. Let $G'_o$ be the projection to $G_o$ of the stabilizer of $o$ in~$G$, that is, the group of all $g_o$ over $g\in G$ fixing $o$. Note that $Y_o$ contains an orbit of~$G_o'$. Thus $G_o'$ has a finite orbit and hence a fixed-point $z_o\in X_o$ by the hypothesis of the lemma. We set $z_p=f(z_o)$ (see Convention \ref{conv:fibre}) for any $f\in G$ with $f(o)=p$, which only depends on $p$ and not on $f$ since $z_o$ was fixed by $G_o'$. }

{Now, let $O\subset P$ be an infinite orbit of $\hbar(G)<H$. Since $Y\subset \Jon$, there is $o\in O$ with $Y_o=\{x_o\}$. Consequently, for any $p\in O$, we have that $Y_p$ consists only of a single element, which we call $z_p$.
Since $Y\subset \Jon$, this $z_p$ equals $x_p$ for all but finitely many $p$.}

{Consequently, we have $z=\{z_p\}_p\in \Jon$. By construction, $z$ is a fixed-point of~$G$.}
\end{proof}

\subsection{Biregularity}

The following encapsulates the idea of a Jonquières transformation $f$ having a persistent base-point in the fiber over a point $p\in P$.

\begin{definition}
\label{def:pers} Let $z=\{z_r\}_r\in \Jon$ be a distinguished vertex. Let $p\in P$ and $f\in G^\oplus$. We say that $f$ is \emph{biregular over $p$} (with respect to $z$) if $f(z)_{f(p)}=z_{f(p)}$ (or, in our notation from Convention~\ref{conv:fibre}, $f(z_p)=z_{f(p)}$). Equivalently, for $f=hg$ with $h\in H, g\in \Pi_pG_0$, we have $g_p(z_p)=z_{f(p)}$ (which equals $z_{h(p)}$). Otherwise, we say that~$f$ is \emph{singular} (with respect to $z$) \emph{over} $p$.

An element $f\in G^\oplus$ \emph{has persistent fiber over $p\in P$} if there exists $l\geq 1$ such that, for all $n\geq l$, we have that $f^n$ is singular over $p$ and $f^{-n}$ is biregular over $p$.
\end{definition}

\begin{remark}
\label{rem:regular}
Note that $f$ is biregular over $p$ if and only if $f^{-1}$ is biregular over~$f(p)$. Furthermore, if $f$ is biregular over $p$ and $f'$ is biregular over $f(p)$, then $f'f$ is biregular over $p$.
\end{remark}

\begin{remark}
\label{rem:fixed}
The point $z$ is a fixed-point for $f$ if and only if $f$ is biregular over all $p\in P$.
\end{remark}

\begin{remark}
\label{rem:per}
If $f$ has persistent fiber over $p$, then it does not have a fixed-point $y\in \Jon$. Indeed, the orbit of $p$ under~$\langle f\rangle$ is infinite, since $f^m(p)=p$ with $m>0$ would imply, by Remark~\ref{rem:regular}, that $f^{l m}$ is simultaneously biregular and singular over~$p$. Furthermore, for all $n\geq l$, we have $z_{f^{n}(p)}\neq f^{n}(z_p)=f^n\big(f^n(z_{f^{-n}(p)})\big)$.
For large~$n$, we have $y_{f^{\pm n}(p)}=z_{f^{\pm n}(p)}$, contradicting $y_{f^{n}(p)}=f^{2n}(y_{f^{-n}(p)})$.
\end{remark}

\subsection{Abelian case}

The following proves Theorem~\ref{thm:maing} in the case where $\hbar(G)$ is abelian. Note that here, as well as in Lemma~\ref{lem:semisimple}, we do not need to assume that $P$ is the projective line $\p^1$.

\begin{lemma}
\label{lem:abel}
Let $G_0$ be a group acting decently on $X_0$. Let $H$ be a group acting on~$P$, and let $G<G^\oplus$ be such that either $\hbar(G)<H$ is trivial or it contains an element that has only finitely many finite orbits on $P$.\footnote{Chris Karpinski proved that this assumption can be removed \cite{Karpinski24}.} If $G$ is finitely generated, acts purely elliptically on $\Jon$, and $\hbar(G)$ is abelian, then $G$ fixes a point of $\Jon$.
\end{lemma}

\begin{proof} If $\hbar(G)$ is trivial, then $G<\Pi_{p\in P} G_p$. Let $Q\subset P$ be a finite set such that each generator $\{g_p\}_{p\in P}$ of $G$ satisfies $g_p(x_p)=x_p$ for $p\notin Q$. For each $q\in Q$, the projection of $G$ to $G_q$ is a purely elliptic subgroup, so it fixes a point $y_q\in X_q$. Setting $y_p=x_p$ for $p\notin Q$, we obtain a fixed-point $\{y_p\}_p$ for $G$ in $\Jon$.

Thus from now on we can assume that there is $t\in G$ with $\hbar(t)$ having only finitely many finite orbits on $P$. Let $z\in \Jon$ be a fixed-point of $t$. Then $t$ is biregular with respect to $z$ over each $p\in P$ {by Remark~\ref{rem:fixed}}.
Let $Q\subset P$ be the union of the finite orbits of~$\langle t\rangle$. Note that $Q$ is preserved by $\hbar(G)$, since $\hbar(G)$ is abelian.

Let $p\in P\setminus Q$. We claim that any $f\in G$ is biregular with respect to $z$ over $p$.
To justify the claim, let $B$ be the finite set of $b\in P\setminus Q$ over which $f$ or $f^{-1}$ is singular. We will show that for some $n>0$, setting $f_n=t^nf$, we have $f_n^i(p)\notin B$, for all $i \in \Z\setminus \{0\}$.
To find such $n$, suppose first that $f^l(p)=t^k(p)$ for some $k,l\in \Z$ with $l> 0$. Then choose $m_0>0$ such that, for all $m\geq m_0$, and all $0\leq j<l$, we~have $t^{\pm m}f^j(p)\notin B$. It suffices then to take $n=m_0+|k|$,
since for any $i\neq 0$ we have $f_n^i(p)=t^{ni}f^i(p)=t^{\pm m}f^j(p)$ for some $0\leq j<l$ and $m\geq m_0$. Otherwise, if there is no $l\neq 0$ with $f^l(p)\in \langle t \rangle(p)$, then
there are finitely many $k\in \Z$ with $t^k(p)\in \langle f\rangle B$, since $B$ is finite. It suffices then to take $n$ larger than the maximum of their $|k|$.

Thus we have $f_n^i(p)\notin B$, for all $i \in \Z\setminus \{0\}$. If $f$ was singular with respect to $z$ over $p$, then by Remark~\ref{rem:regular} $f_n\in G$ would have persistent fiber over $p$ (with $l=1$), contradicting Remark~\ref{rem:per} and justifying the claim.

For each orbit $O$ of $\hbar(G)$ in $Q$, choose $o\in O$ and $y_o\in X_o$ that is fixed by the projection to $G_o$ of the stabilizer of $o$ in $G$, which is of finite index in $G$, hence finitely generated. For each $f\in G$, choose $y_{f(o)}=f(y_o)$ (notation from Convention~\ref{conv:fibre}),
which only depends on $f(o)$ and not on $f$, since $y_o$ was fixed by the projection to $G_o$ of the stabilizer of $o$ in~$G$.
Setting $y_p=z_p$ for $p\notin Q$ gives us a fixed-point $y$ for~$G$.
\end{proof}

\subsection{Semisimple case}

In order to treat the case of non-abelian $\hbar(G)$, we use the dynamics of elements in $\Aut(\p^1)=\PGL_2(k)$.
An element $a\in\PGL_2(k)$ is \emph{semisimple}, if it is conjugate to a diagonal element. With respect to suitable local coordinates, the automorphism of $\p^1$ induced by $a$ is given by $z\mapsto \lambda z$ for some $\lambda\in k^*$. Let us observe that $a$ is of infinite order if and only if $\lambda$ is not a root of unity. In this case, $a$~fixes exactly two points of $\p^1$ and it does not have any other finite orbit. If $k$ admits a norm~$|\cdot|$ with $|\lambda|\neq 1$, then $a$ has north-south dynamics, defined below, in the topology of~$\p^1$ induced by~$|\cdot|$.

\begin{definition}
\label{def:NS}
Let $a$ be a homeomorphism of a Hausdorff topological space $P$, fixing $p,q\in P$, and having the following property. For any disjoint open sets $U\ni p$, $V\ni q$, there is~$n$ such that $a^n(P\setminus V)\subset U$ and $a^{-n}(P\setminus U)\subset V$.
We then say that~$a$ has \emph{north-south dynamics}.
\end{definition}

\begin{lemma}
\label{lem:semisimple}
Let $G_0$ be a group acting decently on $X_0$. Let $H$ be a group acting on $P$. Let $a\in H$ have north-south dynamics {with $p,q$ as in Definition~\ref{def:NS}} for some Hausdorff topology on $P$ in which $H$ acts by homeomorphisms. Suppose also that \hbox{$\mathrm{Stab}(p)\cap\mathrm{Stab}(q)<H$} is abelian.
If $G<G^\oplus$ is finitely generated, acts purely elliptically on~$\Jon$, and $\hbar(G)$ contains $a$, then $G$ has a fixed-point in~$\Jon$.
\end{lemma}

\begin{proof}
Choose $t\in G$ with $\hbar(t)=a$. Let $z\in \Jon$ be a fixed-point of $t$. We claim that for any $f\in G$ and any $r\in P\setminus \{p,q\}$ with $f(r)\neq p,q$, we have that $f$ is biregular with respect to $z$ over $r$.

To justify the claim, first consider the case where $\hbar(f)$ interchanges $p$ and $q$. Then the group $\langle a, \hbar(f)\rangle$ is virtually abelian, {since it has an index $2$ subgroup contained in $\mathrm{Stab}(p)\cap\mathrm{Stab}(q)$}. By Lemma~\ref{lem:abel}, the group $\langle t,f\rangle$ has a finite orbit, hence a fixed-point $y$ in $\Jon$ by Lemma~\ref{lem:first_item}. {Since both $y$ and $z$ belong to $\Jon$}, all but finitely many coordinates of $y$ have to be that of $z$. {Thus} $t(y)=y$ implies $y_r=z_r$ for all $r$ in all infinite orbits of~$\langle t\rangle$, so for all $r\neq p,q$. Since $f$ fixes $y$, we have that $f$ is biregular with respect to $z$ over $r$, as desired.

Second, consider the case where $\hbar(f)$ does not interchange $p$ and $q$. Then after possibly replacing $t$ with $t^{-1}$ and interchanging $p$ with $q$, we can assume $f(p)\neq q$. Let $B\subset P$ be the finite set of points over which $f$ or $f^{-1}$ is singular with respect to $z$. Let $U\ni p$ (respectively, $V\ni q$) be an open neighbourhood intersecting $B\cup \{f^{-1}(p),f(q)\}$ only possibly at $p$ (respectively, $q$), and such that $f(U)$ is disjoint from $V$, which is possible since $f(p)\neq q$.
{Assume also that $U$ is disjoint from $V$.}
Let $n>0$ be as in Definition~\ref{def:NS}.

If $f$ is singular over~$r$, then we have $r,f(r)\in P\setminus V$. Thus $t^nf(r)\subset U$, but $t^nf(r)\neq p$ since $f(r)\neq p$.
Since $f(U)$ is disjoint from $V$, and $U$ does not contain $f^{-1}(p)$, unless $f(p)=p$, we obtain inductively, for all $m>0$, that $(t^nf)^m(r)\in U\setminus p$. We also have $t^{-n}(r)\in V\setminus q$ and analogously we obtain $t^{-n}(t^nf)^m(r)\in V\setminus q$ for all $m<0$. Consequently, $t^nf$ has persistent fiber over $r$ (with $l=1$), which contradicts Remark~\ref{rem:per} and finishes the proof of the claim.

Note that if the entire $\hbar(G)$ stabilizes $\{p,q\}$, then we are done by Lemma~\ref{lem:abel}. Suppose also for the moment that $\hbar(G)$ does not {fix} $p$ or $q$. Then there is $f\in G$ and $r\neq p,q$ with $f(r)=p$. Consider another $f'\in G$ with $r'\neq p,q$ and $f'(r')=p$. Then $f^{-1}f'(r')=r$. Applying the claim above to $f^{-1}f'$, we have $f^{-1}f'(z)_{r'}= z_{r}$. Consequently, $f'(z)_{p}= f(z)_{p}$. We now replace the coordinate $z_p$ of $z$ by $f(z)_{p}$, which, as we have seen, does not depend on~$f$. Note that the new $z$ is still fixed by $t$, which can be verified by substituting above $f'=tf$. Furthermore, now $f$ is biregular over $r$ and $f^{-1}$ is biregular over $p$.
We analogously change the $z_q$ coordinate of~$z$.

We will verify that the new $z$ is a fixed-point for $G$. It remains to verify the biregularity of $f \in G$ over $u\in \{p,q\}$ in the case where $f(u)\in \{p,q\}$. Choose any $r\neq p,q$ and $f'\in G$ with $f'(r)=u$. Introduce $f''=ff'$, which satisfies $f''(r)=f(u)$. From the previous paragraph it follows that both $f',f''$ are biregular over $r$. This implies that $f$ is biregular over $u$, as desired.

In the case where $\hbar(G)$ {fixes}, say, $p$, we redefine $z_p$ to be the fixed-point of the projection of $G$ to $G_p$.
\end{proof}

\subsection{Conclusion}

Recall that an element $h\in\PGL_2(k)$ is \emph{unipotent}, if it is conjugate to an element of the form $\begin{smallpmatrix}
1 & c \\
0 & 1
\end{smallpmatrix}$ for some $c\in k$. By considering the Jordan decomposition, we observe that, for $k$ algebraically closed, every element of $\PGL_2(k)$ is either unipotent or semisimple.

\begin{proof}[Proof of Theorem~\ref{thm:maing}]
Suppose first that $\hbar(G)$ contains a semisimple element $a$ of infinite order with eigenvalues $\lambda, \lambda^{-1}$ and fixed-points $p,q\in \p^1$. Since $a$ is of infinite order, $\lambda$ is not a root of unity. Let $z\in \Jon$. Let $\tilde{k}\subset k$ be the smallest field such that the points $p$ and $q$, the scalar $\lambda$, the elements of $\hbar(G)$, and all the points over which the elements of $G$ are singular with respect to $z$, are defined over $\tilde{k}$. Since $G$ is finitely generated, the field $\tilde{k}$ is a finitely generated field extension over the prime field of~$k$.

Let $\Jon_{\tilde{k}}=\bigoplus_{p\in \p^1(\tilde{k})} (X_p,x_p)$, where $\p^1(\tilde{k})$ is the set of the $\tilde k$-rational points of~$\p^1$. Note that the action of $G$ on $\Jon$ projects to an action on $\Jon_{\tilde{k}}$.

By \cite[Th.\,2.65]{drutukapovich}, we can embed $\tilde{k}$ as a subfield into some local field $K$ with norm $|\cdot|$ such that $|\lambda|\neq 1$. Then $a$ has north-south dynamics on $\p^1(\tilde{k})$ with respect to the Hausdorff topology on $\p^1(\tilde{k})$ induced by the topology of $\p^1(K)$. By Lemma~\ref{lem:semisimple} applied with $P=\p^1(\tilde{k})$, we have that $G$ fixes a point $\{y_p\}_{p\in\p^1(\tilde{k})}\in\Jon_{\tilde{k}}$. Since all the elements of $G$ are biregular with respect to $z$ over all the points outside $\p^1(\tilde{k})$, we~obtain that $G$ fixes the point $\{y_p\}_{p\in \p^1}\in\Jon$, where $y_p=z_p$ for $p\notin\p^1(\tilde{k})$.

Otherwise, if $\hbar(G)$ does not contain a semisimple element of infinite order, then all the elements of $\hbar(G)$ are unipotent or of finite order. By \cite[Prop.\,14.46]{drutukapovich}, there is a finite index subgroup $G'<G$ with $\hbar(G')$ conjugate into the abelian subgroup of the elements of the form $\begin{smallpmatrix}
1 & c \\
0 & 1
\end{smallpmatrix}$.
Lemma~\ref{lem:abel} implies that also in this case $G'$ (and hence~$G$) fixes a point of $\Jon$.
\end{proof}

{\section{Proofs of Theorem~\ref{thm:main_general} and Corollary~\ref{cor:degree}}}

\subsection{Non-rational surfaces}\label{sec:remarks}
Here, we give a proof of Theorem~\ref{thm:main_general}.

\begin{proof}[Proof of Theorem \ref{thm:main_general}]
{Again, by Remark~\ref{rmk_alg_closed_fields}, we can assume that $k$ is algebraically closed}.
If $S$ is rational, then the result follows from Theorem~\ref{thm:main}.

If the Kodaira dimension of $S$ is non-negative, then there exists a smooth projective surface $T$ birationally equivalent to $S$ such that $\Bir(T)=\Aut(T)$ and the result follows from Lemma~\ref{lem:aut}.

Finally, if the Kodaira dimension of $S$ is $-\infty$, but $S$ is not rational, then $S$ is birationally equivalent to $\p^1\times C$ for some non-rational smooth curve $C$. In this case, all the elements $f\in\Bir(S)$ preserve the rational fibration given by the projection to~$C$. {Indeed, if $F$ is a general fiber of the second projection $\pi_2\colon S\to C$, then the restriction of $\pi_2\circ f$ to $F$ induces a rational map to $C$. Since $F\cong \p^1$, this rational map cannot be dominant, hence its image is a point, which implies that $f(F)$ is another fiber.} If the genus of $C$ is $>1$, then $\Bir(C)=\Aut(C)$ is finite. If the genus of~$C$ is~$1$, then $C$ is an elliptic curve and $\Bir(C)=\Aut(C)$ is virtually abelian and virtually consists of translations given by the points $C(\kk)$. We can thus apply Lemma~\ref{lem:abel}.
\end{proof}

\subsection{Degree growth of finitely generated groups}\label{sec:degrees}

There is a well-known and important correspondence between the dynamical behaviour of birational transformations in $\Cr_2(k)$ and the type of isometry they induce on the infinite dimensional hyperbolic space $\h$. The following theorem gives in particular a precise description of the degree growth. It is due to several people. We refer to \cite{cantat2015cremona} for details and references.

\begin{theorem}[Gizatullin; Diller and Favre; Cantat]\label{thm:hyperboloid}
Let $\kk$ be an algebraically closed field and $f\in\Cr_2(k)$. Then one of the following is true:
\begin{enumerate}
\item The transformation $f$ is algebraic, the isometry of $\h$ induced by $f$ is elliptic, and the degree sequence $\{\deg(f^n)\}$ is bounded.
\item[(2a)] The isometry of $\h$ induced by $f$ is parabolic, $\deg(f^n)\sim cn$ for some $c>0$, and $f$ preserves a rational fibration.
\item[(2b)] The isometry of $\h$ induced by $f$ is parabolic, $\deg(f^n)\sim cn^2$ for some $c>0$, and $f$ preserves a fibration of genus $1$ curves.
\item[(3)] The isometry of $\h$ induced by $f$ is loxodromic, $\deg(f^n)=c\lambda^n +\mathcal{O}(1)$ for some $c>0$ and $\lambda>1$.
\end{enumerate}
\end{theorem}

We prove now Corollary~\ref{cor:degree}:

\begin{proof}[Proof of Corollary~\ref{cor:degree}]
If all elements in $G$ induce elliptic isometries on $\h$, \ie they are algebraic, then $G$ and hence $D_T(n)$ are bounded by Theorem~\ref{thm:main}. This implies that the orbit of $G$ on $\h$ is bounded and hence that $G$ fixes a point in~$\h$.

Next, consider the case, where no element in $G$ induces a loxodromic isometry of~$\h$, but there is an element $f\in G$ inducing a parabolic isometry.

First, assume that $f$ preserves a rational fibration. Then all elements in $G$ preserve this same rational fibration (see for instance \cite[Lem.\,5.3.4]{urech2017subgroups}) and after conjugation we may assume that $G$ is a subgroup of the Jonquières group (note that the asymptotic growth of $D_T(n)$ is invariant under conjugation). For an element $g$ in the Jonquières group we have $\deg(g)=\psfrac{\#\bp(g)+1}{2}$ (see for instance \cite{Lamy_Cremona}), where $\#\bp(g)$ denotes the number of base-points of $g$. Since $\#\bp(gh)\leq \#\bp(g)+\#\bp(h)$ we obtain that that $D_T(n)\leq Kn$, where $K=\max_{g\in T}\{\#\bp(g)\}$. At the same time, since by assumption $G$ contains an element whose degree growth is linear, we have $kn\leq D_T(n)$, for some $k>0$. Hence, $D_T(n)\asymp n$.

Now, assume that $f$ preserves a fibration of curves of genus $1$. Again, this implies that all of $G$ preserves the same fibration of curves of genus $1$ and after conjugation we may assume that $G$ is a subgroup of $\Aut(S)$, where $S$ is a Halphen surface \cite[Lem.\,5.3.4]{urech2017subgroups}. In this case, the statement follows from Lemma~\ref{lem:halphen} below.

Finally, we consider the case, where $G$ contains an element $f$ inducing a loxodromic isometry on $\h$. By Theorem~\ref{thm:hyperboloid}, there exist $c$ and $\lambda$ such that $\deg(f^n)=c\lambda^n+\mathcal{O}(1)$. On the other hand, for $\lambda_2=\max_{g\in T}\{\deg(g)\}$ we have $D_T(n)\leq \lambda_2^n$. This shows that $f\asymp \lambda^n$.
\end{proof}

A \emph{Halphen surface} is a rational smooth projective surface $S$ such that $|{-}mK_S|$ is a pencil of genus $1$ curves with empty base locus for some $m>0$. The ideas used in the following lemma have been described in \cite{gizatullin1980rational} (see also \cite{grivaux2013parabolic} and \cite{cantat2012rational}). We~follow the account described in \cite{Lamy_Cremona}.

\begin{lemma}\label{lem:halphen}
Let $S$ be a Halphen surface and let $G<\Aut(S)$ be a finitely generated subgroup containing a non-algebraic element $f$. Then $D_T(n)\asymp cn^2$ for some $c>0$.
\end{lemma}

\begin{proof}
Let us first recall that after possibly passing to a finite index subgroup (which does not change the asymptotic growth of $D_T(n)$), we may assume that $G$ is abelian. Moreover, all algebraic elements in $\Aut(S)$ are of finite order and $\Aut(S)$ does not contain any element inducing a loxodromic isometry on $\h$ (see for instance \cite{cantat2015cremona} for these facts). Hence, again up to passing to a finite index subgroup, we may assume that all elements in $G$ induce a parabolic isometry on~$\h$.

Let $N_{\mathbb{R}}(S)$ be the Néron--Severi space of $S$. Recall that $N_{\mathbb{R}}(S)$ comes with an intersection form of signature $(1,\dim(N_{\mathbb{R}}(S))-1)$, which is preserved by the action of $\Aut(S)$ by push-forwards. There exists a nef divisor class $D_0\in N_{\mathbb{R}}(S)$ such that $D_0\cdot D_0=0$ and such that $g_*D_0=D_0$ for all $g\in G$ (in fact, we can take $D_0=mK_S$). The assumption that all elements in $G$ induce parabolic isometries on $\h$ implies that for all $f\in G$, the only eigenvectors of $f$ are multiples of $D_0$. For all $g\in G$, the restriction of $g_*$ to $D_0^\perp/D_0$ has finite order, since $g_*$ preserves an integral lattice and the induced intersection form on $D_0^\perp/D_0$ is negative definite. Hence, up to passing to a finite index subgroup of $G$, we may assume that the restriction of $G$ to $D_0^\perp/D_0$ is the identity. Let $f_1,\dots, f_k$ be generators of $G$.

Let $A\in N_{\mathbb{R}}(S)$ be an ample divisor. Note that we have $f_*A\neq A$ for all $f\in G$ and $A\notin D_0^\perp$. Write $(f_i)_*A=A+R_i$, where $R_i\in D_0^\perp$. Since the restriction of $(f_i)_*$ to $D_0^\perp/D_0$ is the identity, we can write $(f_i)_*R_j=R_j+t_{ij}D_0$ for some $t_{ij}\in\R$. Let us observe that $(f_i)_*(f_j)_*A=(f_j)_*(f_i)_*A$ for all $i$ and $j$ implies that $t_{ij}=t_{ji}$ for all $i$ and $j$.

By induction, we obtain\vspace*{-3pt}
\[
(f_i)^n_*A=A+nR_i+\frac{n(n-1)t_{ii}}{2}D_0,
\]
and, as a consequence,\vspace*{-3pt}
\[
(f_1)^{n_1}_*\cdots (f_k)^{n_k}_*A=A+\sum_in_iR_i+\biggl(\sum_i\frac{n_i(n_i-1)}{2}t_{ii}+\sum_{i<j}n_in_jt_{ij}\biggr)D_0.
\]

Since, by assumption, all the $f_i$ are non-algebraic and preserving a fibration of genus $1$ curves, the sequence $\deg(f^n)$ grows quadratically in $n$, by Theorem~\ref{thm:hyperboloid}, and hence the $t_{ii}$ are positive. We have that
\[
\deg_A(f_1^{n_1}\cdots f_k^{n_k})=((f_1)^{n_1}_*\cdots (f_k)^{n_k}_*A)\cdot A.
\]
Since $A$ is ample, we have that $A\cdot D_0>0$. Hence, $D_T(n)$ has quadratic growth.
\end{proof}
\Changelback
